% BEGIN VERIFIED DISPLAY
The only unknowns are the thirteen distinct positive integers
\[
J,F,\alpha,f,g,T,X,Y,\xi,\eta,\zeta,\upsilon,\gamma\in\mathbb Z_{>0}.
\]
All other symbols below are polynomial abbreviations:
\[
\begin{gathered}
b=\beta J,\quad q=1+b,\quad C=q-F-\alpha,\quad U=Yq^2,\\
E=XY,\quad k=\eta+\zeta,\quad \ell=EYk,\quad c=\eta+Yk,\\
a=E+Y,\quad H=3+4a,\quad\Delta=H+a^2,\quad L=X+ac,\\
M=q^3b,\quad W=C(X+Kq)-Fq,\\
R=[\alpha+(q-F)b](q^2-1)+J(\mu+\omega q),\\
Z=M(\ell+k+\xi-kW)+Yb-Wq^3,\\
D=UL+H(U\gamma+Eq+R+1-k),\\
I=c^2g,\quad Q=\Delta(f^2-1),\quad V=f^2(Tc-R)-c.
\end{gathered}
\]
The single Diophantine equation is
\[
\boxed{\begin{aligned}
&[1+Z^2-M^2(1+\ell(\ell+k))]\\
&\quad\cdot[1+D^2-U^2(1+\Delta c^2)]\\
&\quad\cdot[\upsilon^2-Q(\upsilon^2-V^2)]
[1+I^2-Q]=1.
\end{aligned}}
\]
% END VERIFIED DISPLAY
